\documentclass[11pt]{article}
\usepackage[utf8]{inputenc}
\usepackage[T1]{fontenc}
\usepackage{amsmath,amssymb,amsthm}
\usepackage{booktabs}
\usepackage{graphicx}
\usepackage{xcolor}
\usepackage{microtype}
\usepackage[margin=1in]{geometry}
\usepackage[hidelinks]{hyperref}
\usepackage{natbib}
\usepackage{enumitem}

\newtheorem{lesson}{Lesson}

\title{\bf Can Bayes Help Retrieve the Right Context for Large Language Model Queries?\\
\large Structure and Adaptation, Not Uncertainty}
\author{Anonymous}
\date{\today}

\begin{document}
\maketitle

\begin{abstract}
Bayesian and uncertainty-quantification language is common in retrieval-augmented generation (RAG) and
text-to-SQL, often without a controlled check that it beats the simple baselines practitioners already
use. We run that check. We treat retrieval as a clean, executable testbed: schema linking for
text-to-SQL, together with multi-hop and single-hop document retrieval. We instrument the five places a
Bayesian idea can enter a retriever, namely calibrated relevance or uncertainty, signal fusion,
structured joint selection, set-level diversity, and online adaptation. Each is compared against a
simple control such as cosine top-$k$, reciprocal-rank fusion, a max-excluded-similarity confidence, or
a fixed top-$k$ context. Across six controlled settings (BIRD, Spider, BEAVER, HotpotQA,
2WikiMultiHopQA, SciFact, and 17 reranking domains) a consistent pattern emerges. Three retrieval ideas
that Bayesian modeling motivates do help. The first is structure, a graph prior over the jointly
relevant subset, which improves multi-hop recall and downstream execution accuracy. The second is
adaptation, online updating, which helps in correlated, repeated workloads. The third is partial
pooling, a hierarchical prior that shares relevance information across domains, which helps in the
scarce-label many-domain regime. Each helps only when its enabling condition holds. A structural prior
is neutral to harmful on single-hop queries, and it can be actively harmful when imposed without genuine
connectivity; adaptation requires a workload that repeats; pooling requires scarce labels. The Bayesian
object tends to fail where it is most often invoked, as an uncertainty or decision signal, losing to
one-line baselines at abstention and context-budget selection, and losing even under distribution shift.
A further finding runs through all three positive cases: the work is done by the modeling idea, not by
the specific Bayesian apparatus. A shortest-path heuristic and a diffusion match the full subgraph
posterior to within a percentage point, a non-Bayesian online counter matches the Bayesian update, and a
frequentist empirical-Bayes matches the hierarchical pooling while a nonparametric escalation over
domains adds nothing. We give four lessons and a practitioner's
decision guide, and we identify the one regime in which the full Bayesian object is worthwhile:
marginalization gates structure across heterogeneous workloads, which provides robustness when one does
not know in advance whether a query needs structure.
\end{abstract}

\section{Introduction}

Retrieval is now a central component of large-language-model (LLM) systems. RAG pipelines
\citep{lewis2020rag} retrieve passages before generating, and text-to-SQL systems retrieve the relevant
tables and columns, a step called schema linking, before writing a query \citep{wang2020ratsql}. In
both, the literature increasingly reaches for Bayesian and uncertainty-quantification (UQ) language:
posterior relevance, calibrated confidence, principled fusion of signals, diversity priors, and
abstention under uncertainty. The literature rarely asks the control question a statistician would ask
first, which is whether any of this improves on the trivial baseline. Cosine top-$k$, reciprocal-rank
fusion, and a single held-out threshold on a transparent score are strong, cheap, and widely used. If a
Bayesian apparatus does not improve on these controls, its practical value is unclear.

This article is that controlled audit. Our answer is conditional, and the conditions are the
contribution. The pattern we find is that Bayesian modeling helps retrieval through the ideas of
structure and adaptation, but not as an uncertainty or decision signal, and that each positive
mechanism helps only under conditions we can identify and map.

\paragraph{What we mean by ``Bayesian.''} To avoid a definitional dispute, we use the term broadly. A
Bayesian idea is a retrieval procedure motivated by a prior, a posterior score, a probabilistic fusion
of signals, a structured joint distribution over the retrieved set, a diversity prior, or online
posterior updating. For each such idea we also evaluate a non-Bayesian control that tests the same
underlying intuition. Examples are a shortest-path graph heuristic against a subgraph posterior, a
max-similarity threshold against a posterior confidence, and an online count against a Bayesian online
update. This separates the value of the idea from the value of the Bayesian machinery, a distinction
that turns out to organize the results.

\paragraph{Why text-to-SQL, and why a multi-setting sweep.} Text-to-SQL retrieval is an unusually clean
testbed. The corpus is a finite, known set of database tables; the gold relevant set is recoverable
from the gold SQL; there is a true relational graph in the foreign keys; and success is checkable by
execution. This lets us measure not only retrieval recall but downstream execution accuracy. We then
test whether each finding recurs in open-domain RAG by replicating on multi-hop QA (HotpotQA,
2WikiMultiHopQA), single-hop retrieval (SciFact), and a correlated enterprise warehouse (BEAVER). The
result is a $5\times5$ map of five Bayesian entry points by five retrieval regimes, summarized in
Figure~\ref{fig:matrix} and Table~\ref{tab:master}.

\paragraph{Why this matters for statisticians.} LLM systems are increasingly built around retrieval,
scoring, calibration, and decision thresholds, all of which are statistical objects. Their Bayesian
interpretations are frequently adopted without the checks a statistician would insist on: whether the
posterior is calibrated, whether it is compared against a sensible baseline, and whether the right loss
is being minimized. This article applies that discipline to a fashionable area. The lessons are
statistical rather than specific to retrieval: a ranking score is not a decision score, an auto-gating
posterior buys robustness rather than accuracy, and a modeling idea can be separated from its Bayesian
apparatus.

\paragraph{Contributions.}
\begin{itemize}[leftmargin=1.4em,itemsep=2pt]
\item A taxonomy of the five places a Bayesian idea enters an LLM retriever, each paired with the
control it must beat (\S\ref{sec:taxonomy}), and a one-page reading guide (\S\ref{sec:guide}).
\item A controlled, reproducible audit across six settings and both domains (\S\ref{sec:results}).
\item The organizing empirical finding, that structure's value tracks the connectivity of the
evidence. It is negative or unnecessary on single-hop queries, it increases with hop depth and corpus
correlation, and it is harmful when an unaligned graph is imposed (\S\ref{sec:boundary}).
\item Evidence that the full Bayesian object is worthwhile in one regime, namely heterogeneous
workloads, where marginalization gates structure automatically (\S\ref{sec:posterior}), and a simple
conceptual account of the results (\S\ref{sec:account}).
\item A constructive method, topology-routed retrieval, and a practitioner's decision guide
(\S\ref{sec:constructive}, \S\ref{sec:guidance}).
\end{itemize}

\section{Background: ``Bayesian retrieval'' already exists}\label{sec:background}

Probabilistic and Bayesian retrieval is not new, and much of what is proposed as ``Bayesian retrieval
for LLMs'' reinvents a mature lineage. The probabilistic relevance framework
\citep{robertson1976relevance} motivates BM25-style scoring \citep{robertson2009bm25} as a
relevance-ranking model, though the resulting scores are not calibrated posterior probabilities.
Language-model IR \citep{ponte1998lm} with Dirichlet smoothing \citep{zhai2004smoothing} is a Bayesian
predictive distribution over terms. Risk-minimization IR \citep{zhai2006risk} casts retrieval as
expected-loss minimization, and diversity via determinantal point processes \citep{kulesza2012dpp} is
an explicitly probabilistic repulsion model. The modern neighborhood our experiments touch includes
dense and late-interaction retrieval \citep{khattab2020colbert}, graph- and multi-hop RAG
\citep{edge2024graphrag,gutierrez2024hipporag}, relation-aware schema linking \citep{wang2020ratsql},
adaptive and active retrieval \citep{jiang2023flare}, and LLM uncertainty and selective prediction
\citep{kuhn2023semantic,wang2023selfconsistency,angelopoulos2021ltt}. Our question is therefore narrow.
Granting that classical probabilistic IR works, does LLM-era Bayes, meaning posteriors over LLM
outputs, graph priors over retrieved sets, and posterior-as-uncertainty, add value beyond cheap
controls, and when? We note where our findings agree with this literature, in that graphs help
multi-hop, and where they sharpen it, by saying when they do not.

\subsection*{What the recent literature believes versus what it shows}

The 2023--2026 literature largely reaches, from different angles, the reading we arrive at by controlled
experiment. There are surveys of uncertainty quantification for large language models
\citep{uqsurvey2025} and a steady stream of ``Bayesian RAG'' and uncertainty-aware retrieval methods,
and the field broadly believes such machinery improves retrieval. What the head-to-head results show is
narrower. The estimators that actually win are consistency- or perturbation-based, sampling several
answers or reasoning paths, or simple discriminative signals, rather than distinctively Bayesian
posteriors. On retrieval-augmented reasoning, white-box token-probability and entropy estimators are the
weakest performers while a consistency estimator leads, and the gain is attributed to explicit
perturbation rather than probability aggregation \citep{soudani2025uqrar}. Semantic entropy does beat
naive token entropy for confabulation detection \citep{farquhar2024semantic}, yet entropy alone admits
confidently-wrong low-entropy errors, so a separate discriminative signal is needed for safe abstention.
Retrieval-augmented model confidence is not robustly calibrated: under irrelevant context it becomes
both over-confident and over-refusing, no single estimator wins across settings, and improving refusal
does not improve calibration \citep{ralm2025know}. Where a probabilistic reranker wins, the gain is an
efficiency trade-off rather than accuracy, and a non-uncertainty static variant already captures most of
it \citep{acurank2025}. Several systems described as probabilistic are not Bayesian at all: conformal RAG
certifies risk with distribution-free frequentist guarantees, and probing-based retrieval decisions are
learned point classifiers. Reports of large gains from ``Bayesian'' retrieval are typically measured
against weak baselines such as BM25 rather than a strong reranker. The consistent picture is that the
field believes Bayes helps, while what is shown is that simple frequentist and consistency baselines are
hard to beat, which is the pattern our audit makes precise.

There is one credible exception, and it is instructive. A query-specific Gaussian-process active-learning
reranker, guided by LLM relevance scores, genuinely beats strong LLM reranking baselines under a matched
query budget \citep{bagel2026}. The value appears exactly where decision theory predicts: when each
relevance judgment is expensive, the useful question is not how to score a fixed set but which candidate
to evaluate next. This is retrieval as Bayesian experimental design, and it is the one regime where our
audit and the literature agree a Bayesian treatment is genuinely worthwhile; we return to it in
\S\ref{sec:limits}.

\section{A taxonomy: five places Bayes can enter an LLM retriever}\label{sec:taxonomy}

We organize the audit by where a Bayesian object could plausibly do work, and pair each with the
control it must beat.

\begin{enumerate}[leftmargin=1.6em,itemsep=3pt]
\item \textbf{Calibrated relevance and uncertainty.} A posterior probability of relevance or answer
correctness, used to rank or to abstain. The control is a held-out threshold on a transparent signal
such as max excluded similarity or a verifier.
\item \textbf{Signal fusion.} A learned or Bayesian combination of heterogeneous scores. The control is
reciprocal-rank fusion (RRF) or plain cosine.
\item \textbf{Structured joint selection.} A prior over the jointly relevant subset, coupling items
through a relational graph of foreign keys, hyperlinks, or entity links. The control is independent
top-$k$, a shortest-path connectivity heuristic, or a diffusion.
\item \textbf{Set-level diversity.} A repulsion prior, here a determinantal point process, to avoid
redundant retrievals. The control is maximal marginal relevance (MMR).
\item \textbf{Online adaptation.} A prior updated from feedback as queries arrive. The control is a
static retriever and a non-Bayesian online counter.
\end{enumerate}

\section{How to read this paper}\label{sec:guide}

Figure~\ref{fig:matrix} gives the whole result before the details. Rows are properties of the retrieval
regime, and columns are the five Bayesian entry points. The summary is as follows. Structure helps when
the evidence is connected, as in multi-hop queries, and hurts when it is not. Diversity helps when
multiple evidence items are independent. Adaptation helps when the workload repeats. Uncertainty and
decision uses of the posterior do not beat simple signals.

\begin{figure}[t]
\centering\small
\renewcommand{\arraystretch}{1.25}
\begin{tabular}{@{}l|ccccc@{}}
\toprule
\textbf{Regime} $\backslash$ \textbf{entry point} & UQ/decision & fusion & structure & diversity & adaptation\\
\midrule
single-item evidence (single-hop)      & ---      & $\approx$ & \textbf{hurts}      & $\approx$            & ---\\
connected evidence (multi-hop)         & ---      & $\approx$ & \textbf{helps}      & neutral             & ---\\
independent multi-evidence             & ---      & $\approx$ & neutral/hurts       & \textbf{helps}      & ---\\
correlated, repeated workload          & ---      & $\approx$ & \textbf{helps more} & ---                 & \textbf{helps}\\
scarce labels, many domains            & loses    & $\approx$ & ---                 & ---                 & \textbf{helps (pool)}\\
heterogeneous (mixed) workload         & loses    & $\approx$ & helps (auto-gated)  & ---                 & ---\\
\bottomrule
\end{tabular}
\caption{Reading guide. ``$\approx$'' marks no improvement over the control, and ``---'' marks an entry
point that is not the operative mechanism in that regime. Bold cells are the main findings. Full numbers
appear in Table~\ref{tab:master}.}
\label{fig:matrix}
\end{figure}

\section{Experimental setup}\label{sec:setup}

\paragraph{Settings and sample sizes.} We span five regimes; Appendix~\ref{app:methods} gives full
detail. S1 is orthogonal SQL, using BIRD \citep{li2023bird} and Spider \citep{yu2018spider}, whose
databases cover unrelated domains. We use an 800-query BIRD sample and the four large-schema BIRD
databases, 436 queries, for downstream execution. S2 is correlated SQL, the BEAVER enterprise warehouse
\citep{chen2024beaver}, with 97 tables and 120 multi-table queries averaging 4 gold tables. S3 is
single-hop: 161 single-table BIRD queries and SciFact \citep{wadden2020scifact}, with 5{,}183 abstracts
and 300 test claims averaging 1.13 gold. S4 is multi-hop RAG, the HotpotQA distractor setting
\citep{yang2018hotpotqa}, with 1{,}492 questions, 10 passages each, 2 gold. S5 is graph RAG,
2WikiMultiHopQA \citep{ho2020twowiki}, with 1{,}500 questions across four reasoning types.

\paragraph{Methods.} For structured selection we build, per query, a cross-fit logistic unary relevance
score over interpretable features: dense cosine, BM25, name and column token overlap, max column
cosine, and value-token overlap. We then couple items through the relational graph in three ways. The
first is an Ising or MRF posterior over item subsets,
$p(x)\propto\exp(\sum_i a_i x_i + \beta\sum_{(i,j)\in E} x_i x_j)$, which does not restrict its support
to connected subsets but favors connectedness through the edge term; we score each item by its
posterior marginal. Marginalization is exact by enumeration over a candidate pool, namely the full
table set for SQL schemas with at most 13 tables and the candidate passages for RAG with at most 16; for
the 97-table BEAVER warehouse we marginalize over the top-15 unary candidates. The second is
personalized PageRank diffusion \citep{page1999pagerank} seeded by the unary scores. The third is a
shortest-path foreign-key closure heuristic. Hyperparameters, namely $\beta$, the diffusion $\alpha$,
and the closure bonus, are selected on held-out queries. Diversity uses MMR \citep{carbonell1998mmr}
and a $k\!=\!2$ DPP. The UQ experiment contrasts the model posterior with a max-excluded-similarity
signal. Adaptation uses online Beta-Bernoulli popularity and a naive-Bayes term-to-table model updated
prequentially, each with a non-Bayesian maximum-likelihood count control. Baselines are cosine top-$k$,
BM25, RRF \citep{cormack2009rrf}, and fixed-$k$ context. We report recall at $|{\rm gold}|$, downstream
execution accuracy (EX) for SQL, and AUROC for abstention; intervals are paired bootstrap unless noted.
Embeddings are \texttt{text-embedding-3-small}, and generation and verification use small chat models.
Total API spend was about \$6. Appendix~\ref{app:methods} gives exact definitions, including
recall@$|{\rm gold}|$, EX, the bridge and comparison labels, graph construction per dataset, and split
logic, and Appendix~\ref{app:sens} gives the sensitivity analyses.

\section{Results, by taxonomy cell}\label{sec:results}

Table~\ref{tab:master} is the main summary. We describe the cells in prose, then devote
\S\ref{sec:boundary} to \S\ref{sec:account} to the findings that organize the rest.

\begin{table}[t]
\centering\small
\caption{Master evidence table. Each row is a contest between a Bayesian object and a control in a given
regime. Numbers are recall at $|{\rm gold}|$ unless noted, and intervals are 95\% paired bootstrap.
Sample sizes are in \S\ref{sec:setup}, and each row maps to a named script in
Appendix~\ref{app:methods}.}
\label{tab:master}
\begin{tabular}{@{}p{3.4cm}p{3.2cm}cp{4.6cm}@{}}
\toprule
\textbf{Cell / regime} & \textbf{Bayesian object vs.\ control} & \textbf{Verdict} & \textbf{Key numbers}\\
\midrule
Structured selection, SQL (S1) & FK-graph MRF vs.\ cosine; vs.\ FK-closure & \textbf{helps} & MRF 0.805/0.822/0.787 vs.\ cosine 0.720/0.673/0.678; vs.\ closure $+$0.02--0.04; \textbf{EX $+$5.7pp} [2.1,9.4]\\
Structured selection, correlated SQL (S2) & join-graph diffusion/MRF vs.\ cosine & \textbf{helps more} & cosine drops to 0.425; PageRank 0.550, MRF 0.556 ($+$12--14pp)\\
Structured selection, multi-hop RAG (S4) & passage-link graph vs.\ cosine & \textbf{helps} & MRF$-$cosine $+$0.104 [.090,.118]; bridge $+$0.17\\
Graph RAG by reasoning type (S5) & entity-graph PageRank/MRF vs.\ cosine & \textbf{helps chained, hurts independent} & bridge $+$0.21, compositional $+$0.17, inference $+$0.13, comparison $-$0.22; PageRank$\approx$MRF (0.805 vs 0.804)\\
Single-hop control, SQL (S3) & FK diffusion/MRF vs.\ cosine & \textbf{structure hurts; MRF neutral} & PageRank 0.727 $<$ cosine 0.839 ($-$0.11); MRF 0.857\\
Single-hop control, RAG (S3) & cosine-kNN graph diffusion vs.\ cosine & \textbf{harmful} & diffusion 0.016 vs.\ cosine 0.608 ($-$0.59)\\
\midrule
Online adaptation, correlated (S2) & online naive-Bayes vs.\ static cosine; vs.\ MLE count & \textbf{helps; idea, not prior} & naive-Bayes $+$0.085 [.062,.113]; MLE count $+$0.090 [.070,.112]; noise- and feedback-robust\\
Few-shot pooling, 17 domains & hierarchical EB vs.\ per-domain/pooled; vs.\ DP-mixture & \textbf{pooling helps; BNP null} & hier beats both extremes each $k$ (CI excludes 0); DP finds 1 cluster; oracle grouping worse than flat\\
\midrule
Decision / context budget (S1) & posterior-threshold context vs.\ fixed-$k$; oracle & \textbf{opportunity real, rule fails} & oracle 0.562@2.2 tbl $>$ full 0.500@10; MRF-threshold 0.411 $<$ fixed-$k$ 0.495\\
Calibrated relevance / UQ & posterior vs.\ cosine max-out & \textbf{loses} & abstention AUROC 0.700 vs.\ 0.763\\
Selective retrieval under shift & deep-ensemble posterior-predictive vs.\ max-softmax & \textbf{loses} & point 0.683$\to$0.755 AUROC under shift; ensemble variance anti-predictive ($<$0.5)\\
Signal fusion (clean text) & learned fusion vs.\ cosine/RRF & \textbf{no improvement} & fusion 0.734 $<$ cosine 0.778\\
Set-level diversity & DPP/MMR vs.\ cosine & \textbf{conditional} & $+$0.07 on comparison, about 0 on bridge; topology-routing $\approx$0.76\\
\bottomrule
\end{tabular}
\end{table}

\paragraph{Structure helps in both domains.} On orthogonal SQL (S1) the FK-graph posterior lifts recall
above cosine and, more importantly, lifts downstream execution accuracy by 5.7 points [2.1, 9.4]. The
mechanism is the recovery of bridge tables, which are join tables whose names share no tokens with the
question and so are invisible to cosine; including them yields correct SQL. The mechanism recurs in
multi-hop RAG (S4), where a passage-link graph beats cosine by 0.104 overall and 0.17 on bridge
questions, and in a second multi-hop dataset (S5).

\paragraph{Adaptation helps in repeated workloads.} Streaming the BEAVER queries and updating an online
term-to-table model prequentially beats static cosine by 8.5 points, with a learning curve in which
recall rises by 0.10 from the first stream quartile to the last. Over 20 random stream orders the gain
is stable, with paired difference 0.085 [.062, .113]. It survives 20\% label noise and degrades
gracefully under partial feedback, falling to a gain of 0.04 at 50\% feedback. A non-Bayesian online
counter without prior smoothing matches the Bayesian version, with gain 0.090 [.070, .112], so the
operative ingredient is online updating rather than the prior.

\paragraph{Partial pooling helps in the few-shot, many-domain regime; nonparametric structure does not.}
The regime that most favors a Bayesian treatment is scarce labels across many related tasks, where
shrinkage should beat a point estimate. We test it with reranking across 17 retrieval domains (five BEIR
datasets and twelve CQADupStack forums), giving each only $k$ labeled queries. A hierarchical
(empirical-Bayes) reranker that partially pools weights across domains beats both a per-domain estimate,
which overfits and falls below the zero-shot cross-encoder at $k\!=\!2$, and a fully pooled estimate, at
every $k$, with the advantage largest at small $k$ and shrinking as labels accumulate: the shrinkage
signature. This is the clearest case in our study of a Bayesian treatment beating a point estimate. Its
distinctively-Bayesian escalation, however, adds nothing. A Dirichlet-process mixture over the domain
weights infers a single cluster, and even an oracle topical grouping of the domains is worse than the
flat hierarchy. Given good shared features, the optimal reranking policy is nearly domain-universal, so
there is no cluster structure to exploit and the nonparametric model correctly collapses to the simple
one. As with structure and adaptation, the useful ingredient is the pooling idea, which a frequentist
empirical-Bayes also captures, not the Bayesian apparatus.

\paragraph{Uncertainty and decision uses lose.} Using the LLM posterior to abstain underperforms a
max-excluded-similarity threshold, with AUROC 0.700 against 0.763. Abstention is scored by AUROC, which
is invariant to any monotone recalibration such as temperature scaling or isotonic regression, so the
gap is a deficit of discrimination rather than calibration, and recalibrating the posterior cannot
close it. Using the subgraph posterior's marginals to choose a variable-size context under-retrieves
and drops execution accuracy to 0.411, below a fixed top-5 at 0.495. The opportunity is real, since the
oracle gold context of 2.2 tables beats the full schema of 10 tables by 6 points, but the raw posterior
scores these models produce are not reliable decision scores without separate validation.

\paragraph{Bayesian uncertainty does not help even under distribution shift.} The usual defense of a
Bayesian treatment is that its value appears out of distribution, where a point confidence is
overconfident. We test this directly. Using a deep ensemble of per-domain reranking heads over the 17
domains above, we predict whether the top retrieved document is relevant (selective retrieval), and
compare a point max-softmax confidence against the ensemble's posterior-predictive uncertainty, in
domain and under a strong shift in which an ensemble built from four source domains is applied to
thirteen distant targets. The point signal improves under shift, from AUROC 0.683 to 0.755, while the
ensemble's epistemic uncertainty is anti-predictive throughout (AUROC below 0.5) and its vote agreement
loses decisively under shift, 0.68 against 0.76. The reason follows from the pooling result: because the
optimal reranking policy is nearly domain-universal, the ensemble members are near-identical, so their
disagreement is noise rather than epistemic signal, and the reranker transfers well enough that shift
does not induce the overconfidence a Bayesian treatment exists to catch.

\paragraph{The UQ failure is not confined to retrieval.} The same posterior-as-uncertainty failure
recurs one layer down, at generation. Used to predict whether a generated SQL query is correct,
black-box LLM-confidence signals reach AUROC at most 0.68, below a separate reasoning verifier at 0.77
and an ensemble at 0.82. For detecting genuinely ambiguous questions, a sampling-posterior divergence
signal is no better than chance, with AUROC about 0.48 and both interpretations jointly covered only
1\% of the time. We report these to show that the pattern is consistent across the stack, covering
ranking, abstention, budgeting, and generation correctness, and we refer the reader to a companion
study for the generation-layer analysis in full.

\paragraph{Fusion and diversity are conditional.} On clean text a learned fusion does not improve on
cosine. Diversity helps only on independent-evidence comparison queries and is neutral on
connected-evidence bridge queries, which motivates the routing result of \S\ref{sec:constructive}.

\section{Structure's value tracks the connectivity of the evidence}\label{sec:boundary}

The main organizing finding is that the sign and size of structure's benefit track how connected the
relevant evidence is. We observe this in every regime we can stratify.

\begin{itemize}[leftmargin=1.4em,itemsep=2pt]
\item \textbf{Single-hop: hard diffusion hurts, marginalization is neutral.} On single-table SQL
queries, PageRank diffusion lowers recall@1 from 0.839 to 0.727, a drop of 0.11, because with one
relevant table diffusion demotes the lone gold toward its graph neighbors. The MRF is neutral to
slightly positive, at 0.857 against cosine at 0.839, because marginalizing over subsets avoids
committing as strongly to propagation. This is a first sign of the auto-gating of \S\ref{sec:posterior}.
On single-hop SciFact the effect of hard diffusion is severe: an imposed cosine-kNN similarity graph
lowers recall from 0.608 to 0.016, because diffusion spreads mass into dense clusters of
topically-similar but irrelevant abstracts. This is a negative control rather than a recommended
SciFact method, and it shows that an arbitrary similarity graph is not a substitute for a
relevance-aligned one. Imposing a graph that is not aligned with relevance can be actively harmful.
\item \textbf{Multi-hop: structure helps, more so with deeper chains.} On 2WikiMultiHopQA the gain
increases with reasoning depth and connectivity: 0.13 for inference, 0.17 for compositional, and 0.21
for bridge-comparison. The one independent-evidence type, comparison, goes the other way, at $-0.22$.
Comparison is a within-dataset negative control, since it uses the same corpus and the same method but
shows the opposite sign, determined by whether the evidence is connected.
\item \textbf{Correlation amplifies the effect.} On the correlated BEAVER warehouse, cosine drops to
0.425, against about 0.72 on orthogonal BIRD, and structure's advantage grows to 12--14 points. The
harder the corpus is for similarity, the more a true relational graph is worth.
\end{itemize}

\paragraph{The graph matters, not the Bayesian machinery.} Wherever structure helps, a personalized
PageRank diffusion and a shortest-path closure match the full subgraph posterior. PageRank and MRF agree
on BIRD, BEAVER, HotpotQA, and 2WikiMultiHopQA. On HotpotQA, where $n=1492$ gives power to detect small
effects, the MRF's edge over PageRank is 0.008 [.0003, .016], which is statistically detectable but
practically negligible (Appendix~\ref{app:sens}). The central ingredient is injecting the correct
relational structure. The embedding-similarity graph is the wrong structure, because connectors are
dissimilar, which is why a correlation prior does not substitute and, on single-hop data, hurts.
Bayesian modeling was useful here for identifying which structure matters, even though the final
practical algorithm can be a one-line heuristic.

\section{When the full posterior is worthwhile}\label{sec:posterior}

There is one regime where the Bayesian object beats the heuristic. On a mixed workload of single- and
multi-hop SQL queries, the subgraph posterior's marginalization gates structure automatically. It
degrades gracefully on single-hop queries, where hard diffusion hurts, and wins on multi-hop queries,
with no predictor. The always-MRF method scores 0.803, against always-PageRank at 0.751 and cosine at
0.744. An explicit hop-count gate cannot recover this, because hop count is hard to predict from the
question, with classifier accuracy 0.799 against 0.936 for the RAG bridge-versus-comparison distinction
of \S\ref{sec:constructive}. The interpretation is statistical rather than about raw accuracy. The full
Bayesian object is not needed to improve accuracy in homogeneous regimes; its value is robustness under
model uncertainty. When one does not know whether a query needs structure, marginalizing over subset
configurations is safer than committing to a hard structural bias.

\section{A conceptual account}\label{sec:account}

The results follow from three simple statements, which we state informally rather than as theorems.

\begin{enumerate}[leftmargin=1.6em,itemsep=2pt]
\item \textbf{Independent retrieval misses bridge evidence.} If a necessary item has low semantic
similarity to the query, such as a join table or a bridge passage, top-$k$ cosine ranks it low and
misses it.
\item \textbf{Graph structure helps only when relevance is connected.} If the gold set is connected in
a relevance-aligned graph, diffusion or subgraph selection propagates score from high-similarity seeds
to the connected bridge items and recovers them. If the gold set is a singleton or its items are
independent, diffusion moves mass away from the answer toward graph neighbors, so the same operation
that helps the connected case hurts the disconnected one. This single mechanism predicts the sign flips
of \S\ref{sec:boundary}.
\item \textbf{A good ranking score need not be a good decision score.} A score that orders items well,
and is therefore useful for retrieval, can still be useless for a threshold decision such as abstention
or budgeting. When the decision is evaluated by a rank-based metric, recalibration cannot help. This
predicts the UQ and context-budget failures.
\end{enumerate}

\section{Routing the prior to the query}\label{sec:constructive}

The audit is not only negative. The right bias is query-dependent, with connectivity for
connected-evidence queries and diversity for independent-evidence queries, so a cheap classifier that
reads the query can route between them. On HotpotQA a lexical logistic classifier predicts bridge
versus comparison at 0.936 accuracy. Routing structure to bridge queries and diversity to comparison
queries reaches recall@2 of 0.748, improving on the best fixed prior at 0.703 by 0.044 [.036, .053] and
nearly matching the oracle at 0.761. The recommendation is not to use a graph prior everywhere but to
match the prior to the query's evidence topology, and, when the topology is hard to classify as in
mixed-hop SQL, to let the posterior's marginalization gate it (\S\ref{sec:posterior}).

\section{Four lessons}\label{sec:lessons}

\begin{lesson}[Structure, not uncertainty]
Bayesian modeling helps when it injects correct structure that a point estimate ignores, here the
foreign-key connectivity that recovers cosine-invisible bridge items. The work is done by the graph,
not by the apparatus: a one-line heuristic matches the posterior, and the embedding-similarity graph,
which is the wrong structure, does not substitute and can hurt.
\end{lesson}

\begin{lesson}[The LLM distribution is a proposal, not a posterior]
LLM output distributions are typically over-confident, so a calibrated decision score cannot be read off
them, and abstention should use a separate transparent signal. Because we evaluate ranking-style
decisions by AUROC, the failure is one of discrimination rather than calibration, and recalibration
cannot fix it. The same miscalibration sinks posterior-as-decision in context budgeting.
\end{lesson}

\begin{lesson}[Adaptation is the other real win, and it too is the idea, not the prior]
Online updating beats a static retriever in correlated, repeated workloads, but a non-Bayesian online
counter matches the Bayesian one, so the value is in updating from feedback rather than in the prior. On
one-shot orthogonal benchmarks there is nothing to adapt to.
\end{lesson}

\begin{lesson}[Marginalization buys robustness, not accuracy]
The one place the full posterior beats the heuristic is heterogeneous workloads, where marginalizing
over subset configurations gates structure on and off without a predictor. Where query types are easy to
classify, a cheap explicit router matches it.
\end{lesson}

\section{Practitioner guidance}\label{sec:guidance}

The following decision procedure summarizes the guidance.
\begin{enumerate}[leftmargin=1.6em,itemsep=1pt]
\item If the answer is supported by a single item, use independent retrieval such as cosine or BM25 and
do not impose a graph.
\item If multiple evidence items are connected by a true relational graph of foreign keys, citations,
hyperlinks, or entity links, use a connectivity method. Start with shortest-path closure or PageRank
diffusion, and reach for the subgraph posterior only on heterogeneous workloads where graceful
auto-gating matters. Expect larger gains as the graph densifies and the corpus correlates.
\item If multiple evidence items are independent facets, use diversity such as MMR or a DPP.
\item If the workload repeats, add online adaptation; a simple online counter suffices.
\item If you need to abstain or budget context, calibrate a held-out threshold on a transparent signal,
and do not threshold the model's posterior.
\end{enumerate}

\section{Scope and limitations}\label{sec:limits}

Our experiments are retrieval-centric and lean on text-to-SQL, where the gold set and the graph are
unambiguous. We treat open-domain RAG through HotpotQA, 2WikiMultiHopQA, and SciFact rather than
web-scale corpora, and we mark SQL-specific claims as such. The patterns we report are consistent across
these controlled settings; we do not claim they exhaust all of IR for LLMs. The adaptation result uses
simulated feedback, with gold revealed after each query, on a modest stream of 120 queries. We report it
with confidence bands over 20 orders and with robustness to partial and noisy feedback, but not at
large-scale online deployment. UQ is evaluated with black-box signals and small models; a dedicated
reasoning verifier does better than both the posterior and the cosine signal, consistent with Lesson 2.
Initial retrieval is dense, via cosine, and lexical, via BM25. We additionally verify that the
structural gain survives a strong cross-encoder reranker (Appendix~\ref{app:sens}), so it is not an
artifact of a weak dense baseline, although we do not exhaustively benchmark every modern retriever.
None of these caveats affects the central, repeatedly observed pattern: the connectivity dependence of
structure and the distinction between the idea and the apparatus.

Finally, one regime we do not study is the strongest remaining case for a Bayesian treatment, and both
our reasoning and the recent literature point to it: budgeted active retrieval. When each relevance
judgment is expensive, for example an LLM call, the operative decision is which candidate to evaluate
next, and treating this as Bayesian experimental design (a posterior over relevance that guides
acquisition) beats fixed-budget reranking under a matched call budget \citep{bagel2026}. This is the
decision layer rather than the ranking or uncertainty layer, and it is where we would expect, and the
evidence supports, that the Bayesian object earns its keep. A controlled audit of budgeted active
retrieval against strong non-Bayesian acquisition rules is the natural sequel to this paper.

\section{Conclusion}

Can Bayesian modeling help retrieval for LLM systems? The practical answer is conditional rather than
uniformly positive. It helps through structure, a graph prior, when evidence is genuinely connected and
the corpus is hard for similarity, and through adaptation, online updating, when the workload repeats.
It does not help as an uncertainty or decision signal, where simple baselines win. Where it does help, a
cheap version of the same idea usually suffices, and the full posterior is reserved for the task it does
best, which is gating structure across heterogeneous queries. For a field that often adopts Bayesian and
uncertainty language around retrieval, the disciplined recommendation is to inject the right structure,
to adapt when the workload allows, and to make decisions with simple, validated signals.

\section*{Reproducibility}
All experiments are scripted and run on cached data. A public repository, with URL withheld for review,
contains the scripts named in Appendix~\ref{app:methods}, an environment file, random seeds, a cached
data manifest, and the exact command to regenerate each table. We use only public benchmarks under their
respective licenses, and we redistribute derived caches rather than raw corpora where licenses preclude
it. API use was safe-by-default, with a dry-run cost estimate, an explicit run flag, and call caps, and
it totaled about \$6.

\appendix
\section{Detailed methods and definitions}\label{app:methods}

\paragraph{Recall at $|{\rm gold}|$.} For a query with gold set $G$, retrieve the top $|G|$ items under
the method's score and report $|G \cap \widehat{G}|/|G|$, averaged over queries. For single-hop queries,
where $|G|\!=\!1$, this is recall@1.

\paragraph{Downstream execution accuracy.} For the four large-schema BIRD databases, we prune the schema
shown to the generator to a method's selected tables, generate SQL with a small chat model, and execute
against the full database. EX is the fraction whose result set matches the gold query's, guarded by an
execution timeout.

\paragraph{Evidence types.} Bridge, compositional, and inference questions require chained evidence,
where item B is found via item A. Comparison questions require independent evidence about two entities.
For HotpotQA we use the provided type label, and for 2WikiMultiHopQA the four native types. A query is
connected if its gold items form a connected subgraph of the relevance graph.

\paragraph{Graph construction.} For SQL, nodes are tables and edges are foreign keys; for BEAVER the
declared join keys are collapsed to table-table edges. For HotpotQA and 2WikiMultiHopQA, nodes are
candidate passages, with a directed edge from $i$ to $j$ when passage $j$'s title is mentioned in
passage $i$, a title-mention link graph. SciFact has no native graph, so we construct a cosine-$k$NN
graph with $k\!=\!8$ purely to test whether an imposed similarity graph helps single-hop retrieval,
which it does not.

\paragraph{Unary score and splits.} The unary relevance model is logistic regression on the features in
\S\ref{sec:setup}, fit with two-fold cross-fitting at the query level, so that no item from a test query
informs its own score. Hyperparameters, namely $\beta$, $\alpha$, and the closure bonus, are chosen on
the training fold. Bootstrap CIs resample queries, paired across methods, 2{,}000 to 3{,}000 times.

\paragraph{Scripts.} \texttt{s1a\_graphgp.py} for S1, comparing diffusion and MRF; \texttt{s2\_beaver.py}
for S2; \texttt{s3\_sql\_singlehop.py} and \texttt{s3\_sql\_hopgate.py} for single-hop SQL and the
mixed-workload auto-gating result; \texttt{s3\_scifact.py} for single-hop RAG; \texttt{s4\_hotpot.py}
and \texttt{s4\_diversity\_uq.py} for S4 structure, diversity, and UQ; \texttt{s4\_adaptive.py} for
topology routing; \texttt{s5\_twowiki.py} for S5; \texttt{s\_adapt\_beaver.py} for adaptation;
\texttt{downstream\_ex.py} for EX and context budget; and \texttt{s\_sensitivity.py} for
Appendix~\ref{app:sens}.

\section{Sensitivity and robustness}\label{app:sens}

\textbf{Diffusion strength.} On HotpotQA, recall@2 is stable across the PageRank restart parameter,
ranging from 0.710 at $\alpha\!=\!0.3$ to 0.703 at $\alpha\!=\!0.85$.

\textbf{Retrieval depth.} Structure helps at $k\!\ge\!2$. At recall@3, cosine is 0.780, PageRank 0.803,
and MRF 0.808. It slightly hurts at $k\!=\!1$, where cosine is 0.435 and PageRank 0.397, consistent with
the connectivity account.

\textbf{MRF versus heuristic.} The paired MRF$-$PageRank difference in recall@2 is 0.008 [.0003, .016]
at $n\!=\!1492$, which is detectable but about one point.

\textbf{Strong reranker.} The structural gain is not an artifact of a weak base ranker. Replacing the
cosine seed with an ms-marco cross-encoder, \texttt{ms-marco-MiniLM-L-6-v2}, on HotpotQA, the
cross-encoder alone reaches recall@2 of 0.720, against cosine at 0.685. Adding PageRank structure on top
raises bridge-question recall by 0.067 [.050, .084] and overall recall by 0.023 [.006, .039], and the
MRF again degrades more gracefully on comparison questions, at 0.822 against PageRank at 0.700. Structure
is complementary to deep cross-attention rather than subsumed by it.

\textbf{Adaptation.} Over 20 stream orders the online gain is 0.085 [.062, .113] for the Bayesian model
and 0.090 [.070, .112] for the non-Bayesian MLE count. Under 20\% label noise the all-stream recall is
unchanged at 0.510, and under 50\% and 25\% feedback it falls gracefully to 0.469 and 0.437.

\bibliographystyle{plainnat}
\bibliography{tas_refs}

\end{document}
